% ============================================================
\subsection{Propagative Persistence and Reversible Wave Realization}
\label{sec:flagship-wave}
% ============================================================

The wave lane is the realization of \emph{propagative persistence}: scalar
coherence remains visible by reversible spread through the realized temporal
order. The physical question on this lane is therefore how persistent scalar
coherence propagates on the time-realizing cut. Second-order transport
rigidity, scalar endpoint closure, and the continuum bridge then determine the
physical wave law class using only the machinery inherited from the earlier
structural chapters together with the continuum bridge.
The canonical wave form is obtained by identification of the local endpoint
representative inside that forced class. The exact
constant-mass Klein--Gordon specialization is a further refinement after that
local collapse, not part of the forcing of the wave representative itself.

Theorem~\ref{thm:temporal-realization-forbids-primitive-internal-creation}
again rules out primitive source creation on a closed reversible carrier. Any
visible inhomogeneity in the selected wave law must therefore be an exchange
remainder generated by coupling or hidden-sector return.

The relevant compatibility identity on this carrier is reversible second-order
transport: the selected cut is the smallest one on which the scalar wave
transport datum and scalar endpoint closure class are faithfully visible. The
interval \(J\subseteq\bbR\) introduced below is therefore a coordinate model of
the realized causal order on this carrier, not an extra primitive input.

\paragraph{Carrier status.}
The wave carrier itself is already fixed at the common-law level. By
Theorem~\ref{thm:three-forcing-bridges-from-bedrock-to-primitive-physics}, the
primitive direct lanes are already forced on the selected bundle, and the wave
lane is the canonical realization of the \(2\)-channel. The realization and
endpoint packages below are therefore compatibility/readout data for an
already forced lane, not an independent wave-sector axiom.

\subsubsection{Reversible Wave Realization}

\begin{definition}[Reversible wave realization]
\label{def:reversible-wave-realization}\lean{WaveLane.ReversibleRealization}
Let \(J\subseteq \bbR\) be a compact time interval, let \(\mathcal W\) be a
Hilbert space with dense test domain \(\mathcal W_0\subseteq\mathcal W\), and
let
\[
\mathcal F=\{(R_n,h_n,L_n)\}_{n\in\bbN}
\]
be a refinement-compatible realized law family into \(\mathcal W\). A
\emph{reversible wave realization} on a visible carrier
\(\mathcal X_{\mathrm{wav}}\subseteq\mathcal W\) consists of:
\begin{enumerate}[label=(\roman*),itemsep=0.2em]
\item for every \(n\), the relevant selected channels of \(R_n\) carry
      selected channel transport systems and belong to admissible realized
      classes satisfying the hypotheses of
      Corollary~\ref{cor:intrinsic-package-implies-governing-datum-rigid};
\item the selected channel families are organized into a
      refinement-compatible family on \(C^2(J,\mathcal W_0)\) with a
      refinement-coherent reconstruction datum whose exact core realization is
      bounded;
\item the induced continuum transport class on \(\mathcal X_{\mathrm{wav}}\)
      is represented by a self-adjoint nonnegative second-order operator
      \[
      \mathcal T_{\chi}:\mathcal W_0\to\mathcal W;
      \]
\item the induced continuum closure class on \(\mathcal X_{\mathrm{wav}}\) is
      represented by a self-adjoint scalar operator
      \[
      M_{\chi}:\mathcal W_0\to\mathcal W.
      \]
\end{enumerate}
\end{definition}

\subsubsection{Least-Motion Wave Carrier}

\begin{theorem}[Least-motion wave observer]
\label{thm:least-motion-wave-observer}\lean{WaveLane.leastMotionObserver}
Assume the hypotheses of Definition~\ref{def:reversible-wave-realization}.
Suppose:
\begin{enumerate}[label=(\roman*),itemsep=0.2em]
\item the wave carrier \(\mathcal X_{\mathrm{wav}}\) is faithful,
      closure-admissible, and stable under admissible promotion for the
      realized family;
\item no proper visible subcarrier of \(\mathcal X_{\mathrm{wav}}\) still
      supports a forced scalar second-order transport datum together with a
      scalar endpoint closure class at the observed scale;
\item any characteristic observer datum satisfying \emph{(i)}--\emph{(ii)} is
      horizon-preserving equivalent to the datum generated by
      \(\mathcal X_{\mathrm{wav}}\).
\end{enumerate}
Then \(\mathcal X_{\mathrm{wav}}\) is the unique least-motion characteristic
observer for the realized family.
\end{theorem}

(Proof in Appendix~\ref{sec:appendix-carrier}.)

\subsubsection{Induced Reversible Wave Law}

\begin{theorem}[Induced reversible wave law on the selected cut]
\label{thm:induced-wave-law}\lean{WaveLane.inducedLaw}
Assume the hypotheses of
Theorem~\ref{thm:least-motion-wave-observer}. Then there exist a unique induced
speed coefficient \(c_{\chi}>0\), a self-adjoint nonnegative transport operator
\(\mathcal T_{\chi}\), a self-adjoint scalar closure operator \(M_{\chi}\), and
a remainder \(\mathsf R_{\mathrm{wav}}\in C(J,\mathcal W)\) such that every
\(\phi\in C^2(J,\mathcal W_0)\) satisfies
\[
\partial_t^2\phi
+
c_{\chi}^2\mathcal T_{\chi}\phi
+
M_{\chi}\phi
=
\mathsf R_{\mathrm{wav}}.
\]
By
Theorem~\ref{thm:temporal-realization-forbids-primitive-internal-creation},
\(\mathsf R_{\mathrm{wav}}\) is an exchange remainder rather than an
independent source term.
If the selected cut is lossless at leading order, so that
\(\mathsf R_{\mathrm{wav}}=0\), this is the forced reversible wave law on the
selected cut.
\end{theorem}

(Proof in Appendix~\ref{sec:appendix-carrier}.)

The theorem above therefore reduces the wave lane to a local
representative-identification problem. The selected cut already fixes the
second-order transport operator, the scalar closure operator, and the
admissible exchange remainder. The canonical pairing-symmetric representative
of this forced class is then identified with the standard
wave/Laplace/scalar-profile operator. Once that identification is made, the
recognizable scalar wave identity follows; the exact Klein--Gordon
specialization is then the additional statement that the forced scalar profile
has a distinguished constant-mass part.

\subsubsection{Endpoint Identification and Wave Prediction}

In standard relativistic scalar notation, the recognized endpoint law is the
scalar wave equation with real scalar profile. Here again the final
identification is not a free notation choice: it is the local endpoint collapse
of the already forced wave class. The nearest-neighbor specialization below is
then the standard periodic-lattice dispersion law
\cite{Klein1926,Gordon1926,Brillouin1953}.

\begin{definition}[Wave local endpoint-rigidity package]
\label{def:wave-local-endpoint-rigidity-package}\lean{WaveLane.EndpointCollapse}
Assume the hypotheses of Theorem~\ref{thm:induced-wave-law}. A
\emph{wave local endpoint-rigidity package} on the selected wave carrier
\(\mathcal X_{\mathrm{wav}}\) consists of:
\begin{enumerate}[label=(\roman*),itemsep=0.2em]
\item the minimal pairing-symmetric transport family on
      \(\mathcal X_{\mathrm{wav}}\) is represented levelwise by a
      refinement-compatible local scalar stencil family of minimal local order
      \(2\);
\item on each relevant selected wave channel, the corresponding one-step
      transport data are without preferred direction at quadratic order in the
      sense of
      Definition~\ref{def:no-preferred-direction-selected-channel-system};
\item the transport stencil family carries a refinement-coherent
      reconstruction datum on \((\mathcal W,\mathcal W_0)\) whose exact core
      realization is bounded;
\item the scalar endpoint-closure family on \(\mathcal X_{\mathrm{wav}}\) is
      represented levelwise by a refinement-compatible local scalar stencil
      family of order \(0\);
\item the scalar closure family carries a refinement-coherent reconstruction
      datum on \((\mathcal W,\mathcal W_0)\) whose exact core realization is
      bounded.
\end{enumerate}
\end{definition}

\begin{theorem}[Wave local endpoint collapse to Laplace/scalar-profile form]
\label{thm:wave-local-endpoint-collapse}\lean{WaveLane.endpointCollapse}
Assume the hypotheses of
Definition~\ref{def:wave-local-endpoint-rigidity-package}. Then there exist a
forced positive Laplace-type representative
\[
\Delta_{\chi}:\mathcal W_0\to\mathcal W
\]
for the selected transport class and a forced real multiplication operator
\[
W_{\chi}:\mathcal W_0\to\mathcal W
\]
for the selected scalar closure class such that every
\(\phi\in C^2(J,\mathcal W_0)\) satisfies
\[
\partial_t^2\phi
-
c_{\chi}^2\Delta_{\chi}\phi
+
W_{\chi}\phi
=
\mathsf R_{\mathrm{wav}}.
\]
If the selected cut is lossless at leading order, so that
\(\mathsf R_{\mathrm{wav}}=0\), this is exactly the scalar wave equation on the
selected wave carrier.
\end{theorem}

(Proof in Appendix~\ref{sec:appendix-carrier}.)

\begin{corollary}[Recognizable wave/Klein--Gordon identity]
\label{cor:recognizable-wave-kg-identity}\lean{WaveLane.recognizableIdentity}
Assume the hypotheses of
Theorem~\ref{thm:wave-local-endpoint-collapse}. Assume in addition that the
forced real scalar closure profile on the selected carrier is written in the
form
\[
W_{\chi}=m_{\chi}^2+V
\]
for a nonnegative constant mass datum \(m_{\chi}^2\) and a real residual
potential \(V\). Then
\[
\partial_t^2\phi
-
c_{\chi}^2\Delta_{\chi}\phi
+
(m_{\chi}^2+V)\phi
=
\mathsf R_{\mathrm{wav}}.
\]
If the selected cut is lossless at leading order, so that
\(\mathsf R_{\mathrm{wav}}=0\), this is exactly the scalar wave identity on the
selected carrier, and when \(m_{\chi}\neq 0\) it is its Klein--Gordon form.
\end{corollary}

\begin{proof}
This is exactly Theorem~\ref{thm:wave-local-endpoint-collapse} after writing the
forced real multiplication profile \(W_{\chi}\) in the displayed
mass/potential form.
\end{proof}

\begin{corollary}[Exact Schur form of the selected wave remainder]
\label{cor:wave-schur-remainder-identity}\lean{WaveLane.schurRemainder}
Assume the hypotheses of Theorem~\ref{thm:wave-local-endpoint-collapse}. Assume
in addition that the same selected datum satisfies a canonical selected Schur
bridge. Then there is a unique wave-projected Schur correction
\[
\Sigma_{h,\mathrm{wav}}(z)\in C(J,\mathcal W)
\]
obtained by reading the common selected Schur correction \(\CCSigma{h}{z}\)
through the forced wave endpoint identification on the selected carrier, and
\[
\mathsf R_{\mathrm{wav}}=\Sigma_{h,\mathrm{wav}}(z).
\]
Accordingly, every \(\phi\in C^2(J,\mathcal W_0)\) satisfies
\[
\partial_t^2\phi
-
c_{\chi}^2\Delta_{\chi}\phi
+
W_{\chi}\phi
=
\Sigma_{h,\mathrm{wav}}(z).
\]
\end{corollary}

(Proof in Appendix~\ref{sec:appendix-carrier}.)

\begin{definition}[Homogeneous constant-mass wave refinement package]
\label{def:homogeneous-constant-mass-wave-refinement-package}\lean{WaveLane.ConstantMassRefinement}
Assume the hypotheses of
Theorem~\ref{thm:wave-local-endpoint-collapse}. A
\emph{homogeneous constant-mass wave refinement package} on the selected wave
carrier consists of:
\begin{enumerate}[label=(\roman*),itemsep=0.2em]
\item on each relevant selected wave channel, the admissible realized class
      furnishing the scalar closure profile is locally homogeneous on the
      selected channels in the sense of
      Definition~\ref{def:locally-homogeneous-realized-class};
\item the selected scalar closure family is equivariant for the induced local
      symmetry action on the corresponding homogeneous selected-wave sector;
\item after the wave local endpoint collapse, the forced zeroth-order closure
      profile on the selected wave carrier has no residual site-dependent
      scalar defect beyond that homogeneous selected-wave sector;
\item the resulting scalar closure is nonnegative on that homogeneous
      selected-wave sector.
\end{enumerate}
\end{definition}

\begin{theorem}[Homogeneous wave scalar profile forces exact constant mass]
\label{thm:homogeneous-wave-scalar-profile-forces-constant-mass}\lean{WaveLane.constantMass}
Assume the hypotheses of
Definition~\ref{def:homogeneous-constant-mass-wave-refinement-package}. Then
there exists a unique nonnegative constant mass datum \(m_{\chi}^2\in\bbR\)
such that
\[
W_{\chi}=m_{\chi}^2 I
\]
on the selected wave carrier. Consequently every
\(\phi\in C^2(J,\mathcal W_0)\) satisfies
\[
\partial_t^2\phi
-
c_{\chi}^2\Delta_{\chi}\phi
+
m_{\chi}^2\phi
=
\mathsf R_{\mathrm{wav}}.
\]
\end{theorem}

(Proof in Appendix~\ref{sec:appendix-carrier}.)

\begin{corollary}[Exact recognizable Klein--Gordon identity]
\label{cor:exact-recognizable-klein-gordon-identity}\lean{WaveLane.exactKleinGordon}
Assume the hypotheses of
Theorem~\ref{thm:homogeneous-wave-scalar-profile-forces-constant-mass}. Then
\[
\partial_t^2\phi
-
c_{\chi}^2\Delta_{\chi}\phi
+
m_{\chi}^2\phi
=
\mathsf R_{\mathrm{wav}}.
\]
If the selected cut is lossless at leading order, so that
\(\mathsf R_{\mathrm{wav}}=0\), this is exactly the constant-mass
Klein--Gordon equation on the selected wave carrier.
\end{corollary}

\begin{proof}
This is exactly
Theorem~\ref{thm:homogeneous-wave-scalar-profile-forces-constant-mass} read in
the standard Klein--Gordon notation.
\end{proof}

\begin{definition}[Nearest-neighbor wave specialization]
\label{def:nearest-neighbor-wave-specialization}\lean{WaveLane.NearestNeighborSpecialization}
Assume the hypotheses of
Theorem~\ref{thm:induced-wave-law}. A \emph{nearest-neighbor wave
specialization} of the selected cut consists of:
\begin{enumerate}[label=(\roman*),itemsep=0.2em]
\item a one-dimensional primitive lattice with spacing \(\ell_{\chi}>0\);
\item the lossless massless leading-order regime
      \[
      M_{\chi}=0,
      \qquad
      \mathsf R_{\mathrm{wav}}=0;
      \]
\item the selected transport operator acting by the second-difference law
      \[
      (\mathcal T_{\chi}\phi)_n
      =
      \frac{2\phi_n-\phi_{n+1}-\phi_{n-1}}{\ell_{\chi}^2}.
      \]
\end{enumerate}
\end{definition}

\begin{theorem}[Exact lattice dispersion on the selected wave carrier]
\label{thm:exact-wave-dispersion}\lean{WaveLane.exactDispersion}
Assume the hypotheses of
Definition~\ref{def:nearest-neighbor-wave-specialization}. Then plane-wave
modes
\[
\phi_n(t)=e^{i(nk\ell_{\chi}-\omega t)}
\]
satisfy
\[
\omega(k)^2
=
\frac{4c_{\chi}^2}{\ell_{\chi}^2}
\sin^2\!\left(\frac{k\ell_{\chi}}{2}\right),
\]
and, writing \(E=\hbar\omega\) and \(E_{\chi}:=\hbar c_{\chi}/\ell_{\chi}\),
the corresponding group velocity is
\[
v_g(E)
=
c_{\chi}\sqrt{1-\left(\frac{E}{2E_{\chi}}\right)^2}
\qquad (0\le E\le 2E_{\chi}).
\]
\end{theorem}

(Proof in Appendix~\ref{sec:appendix-carrier}.)

\begin{corollary}[Observable wave-carrier signature]
\label{cor:observable-wave-carrier-signature}\lean{WaveLane.observableCarrierSignature}
Assume the hypotheses of
Theorem~\ref{thm:exact-wave-dispersion}. Then in the low-energy regime
\(E\ll E_{\chi}\),
\[
\frac{\delta v}{c_{\chi}}
=
\frac{v_g(E)-c_{\chi}}{c_{\chi}}
\approx
-\frac18\left(\frac{E}{E_{\chi}}\right)^2.
\]
In particular, the sign is fixed and subluminal, and the selected wave carrier
also predicts a hard cutoff at \(E=2E_{\chi}\), where \(v_g(E)\) collapses to
zero.
\end{corollary}

(Proof in Appendix~\ref{sec:appendix-carrier}.)

\subsubsection{Schur-Corrected Lattice Dispersion}

The lossless nearest-neighbor law isolates the clean leading wave surface.
To read the predictive remainder without weakening that identification, keep
the same lattice transport and constant-mass closure but allow the exact
returned Schur correction to survive on a single selected plane-wave sector.
On such a sector the returned wave remainder is read as a scalar correction to
the dispersion relation, not as a new primitive law.

\begin{definition}[Mode-compatible Schur-corrected lattice-wave specialization]
\label{def:mode-compatible-schur-corrected-lattice-wave-specialization}\lean{WaveLane.CorrectedDispersion.LatticeSpecialization}
Assume the hypotheses of
Corollary~\ref{cor:wave-schur-remainder-identity} and
Theorem~\ref{thm:homogeneous-wave-scalar-profile-forces-constant-mass}. A
\emph{mode-compatible Schur-corrected lattice-wave specialization} of the
selected cut consists of:
\begin{enumerate}[label=(\roman*),itemsep=0.2em]
\item a one-dimensional primitive lattice with spacing \(\ell_{\chi}>0\);
\item the selected Laplace-type transport representative acting by the
      second-difference law
      \[
      (\Delta_{\chi}\phi)_n
      =
      \frac{\phi_{n+1}-2\phi_n+\phi_{n-1}}{\ell_{\chi}^2};
      \]
\item the selected scalar closure being the constant-mass operator
      \[
      W_{\chi}=m_{\chi}^2 I
      \qquad (m_{\chi}\ge 0);
      \]
\item a selected plane-wave sector generated by
      \[
      \phi_n(t)=e^{i(nk\ell_{\chi}-\omega t)};
      \]
\item the exact returned wave remainder being mode-compatible and time-local on
      that sector, in the sense that there exists a real scalar correction
      profile \(\sigma_{\chi}(k;z)\) such that
      \[
      \Sigma_{h,\mathrm{wav}}(z)_n(t)
      =
      \sigma_{\chi}(k;z)\,\phi_n(t)
      \]
      for every mode in the selected sector.
\end{enumerate}
\end{definition}

\begin{theorem}[Exact Schur-corrected lattice dispersion on the selected wave carrier]
\label{thm:exact-schur-corrected-wave-dispersion}\lean{WaveLane.CorrectedDispersion.laws}
Assume the hypotheses of
Definition~\ref{def:mode-compatible-schur-corrected-lattice-wave-specialization}.
Then the selected plane-wave sector satisfies the exact corrected dispersion law
\[
\omega(k;z)^2
=
m_{\chi}^2
+
\frac{4c_{\chi}^2}{\ell_{\chi}^2}
\sin^2\!\left(\frac{k\ell_{\chi}}{2}\right)
-
\sigma_{\chi}(k;z).
\]
If in addition \(k\mapsto \sigma_{\chi}(k;z)\) is differentiable and the
displayed right-hand side is positive, then the corresponding group velocity is
\[
v_g(k;z)
=
\frac{
\frac{2c_{\chi}^2}{\ell_{\chi}}\sin(k\ell_{\chi})
-\partial_k\sigma_{\chi}(k;z)
}{
2\sqrt{
m_{\chi}^2
+
\frac{4c_{\chi}^2}{\ell_{\chi}^2}
\sin^2\!\left(\frac{k\ell_{\chi}}{2}\right)
-
\sigma_{\chi}(k;z)
}
}.
\]
When \(\sigma_{\chi}(k;z)=0\), this reduces to the lossless lattice law of
Theorem~\ref{thm:exact-wave-dispersion}; when \(m_{\chi}=0\), it further
reduces to the massless signature of
Corollary~\ref{cor:observable-wave-carrier-signature}.
\end{theorem}

(Proof in Appendix~\ref{sec:appendix-carrier}.)

This theorem gives an exact non-lossless wave specialization on the selected
plane-wave sector. Once a concrete realized family yields a returned profile
\(\sigma_{\chi}(k;z)\), the observed dispersion and group velocity are fixed
by the displayed formulas. The lossless lattice law is recovered as the
special case \(\sigma_{\chi}(k;z)=0\).

\subsubsection{Concrete Hidden-Channel Model}

The abstract corrected profile \(\sigma_{\chi}(k;z)\) becomes more informative
once one computes it on a concrete realized family. The simplest exactly
solvable model is a visible lattice wave locally coupled to one hidden lattice
wave on the same selected primitive lattice. In that case the returned
remainder is a genuine wave-side self-energy induced by exact elimination of
the hidden channel.

\begin{definition}[Coupled visible-shadow lattice-wave model]
\label{def:coupled-visible-shadow-lattice-wave-model}\lean{WaveLane.HiddenChannel.LatticeModel}
Assume the hypotheses of
Definition~\ref{def:mode-compatible-schur-corrected-lattice-wave-specialization}.
A \emph{coupled visible-shadow lattice-wave model} consists of:
\begin{enumerate}[label=(\roman*),itemsep=0.2em]
\item a hidden scalar lattice field \(\eta_n(t)\) on the same one-dimensional
      lattice of spacing \(\ell_{\chi}\);
\item a hidden nearest-neighbor wave operator with parameters
      \(\widetilde c_{\chi}>0\) and \(\widetilde m_{\chi}\ge 0\), acting by
      \[
      \partial_t^2\eta
      -
      \widetilde c_{\chi}^2\Delta_{\chi}\eta
      +
      \widetilde m_{\chi}^2\eta;
      \]
\item a real local coupling amplitude \(\gamma_{\chi}\ge 0\) such that the
      visible and hidden fields satisfy
      \[
      \partial_t^2\phi
      -
      c_{\chi}^2\Delta_{\chi}\phi
      +
      m_{\chi}^2\phi
      =
      \gamma_{\chi}\eta,
      \]
      \[
      \partial_t^2\eta
      -
      \widetilde c_{\chi}^2\Delta_{\chi}\eta
      +
      \widetilde m_{\chi}^2\eta
      =
      \gamma_{\chi}\phi;
      \]
\item a common plane-wave sector on which
      \[
      \phi_n(t)=\Phi\,e^{i(nk\ell_{\chi}-\omega t)},
      \qquad
      \eta_n(t)=H\,e^{i(nk\ell_{\chi}-\omega t)}
      \]
      for amplitudes \(\Phi,H\in\mathbb C\).
\end{enumerate}
\end{definition}

\begin{theorem}[Exact hidden-channel wave remainder and branch splitting]
\label{thm:hidden-channel-wave-remainder-and-branch-splitting}\lean{WaveLane.HiddenChannel.splitting}
Assume the hypotheses of
Definition~\ref{def:coupled-visible-shadow-lattice-wave-model}. Define the
visible and hidden leading profiles
\[
\omega_{\mathrm{vis}}(k)^2
:=
m_{\chi}^2
+
\frac{4c_{\chi}^2}{\ell_{\chi}^2}
\sin^2\!\left(\frac{k\ell_{\chi}}{2}\right),
\]
\[
\omega_{\mathrm{hid}}(k)^2
:=
\widetilde m_{\chi}^2
+
\frac{4\widetilde c_{\chi}^2}{\ell_{\chi}^2}
\sin^2\!\left(\frac{k\ell_{\chi}}{2}\right).
\]
Then on every mode with \(\omega_{\mathrm{hid}}(k)^2\neq \omega^2\), exact
elimination of the hidden amplitude gives the returned wave profile
\[
\sigma_{\chi}^{\mathrm{hid}}(k,\omega)
=
\frac{\gamma_{\chi}^2}{\omega_{\mathrm{hid}}(k)^2-\omega^2}.
\]
Accordingly the visible dispersion law is
\[
\omega^2
=
\omega_{\mathrm{vis}}(k)^2
-
\sigma_{\chi}^{\mathrm{hid}}(k,\omega),
\]
equivalently
\[
(\omega^2-\omega_{\mathrm{vis}}(k)^2)
(\omega^2-\omega_{\mathrm{hid}}(k)^2)
=
\gamma_{\chi}^2.
\]
Hence the two exact propagating branches are
\[
\omega_{\pm}(k)^2
=
\frac{
\omega_{\mathrm{vis}}(k)^2+\omega_{\mathrm{hid}}(k)^2
\pm
\sqrt{
\bigl(\omega_{\mathrm{vis}}(k)^2-\omega_{\mathrm{hid}}(k)^2\bigr)^2
+
4\gamma_{\chi}^2
}
}{2}.
\]
\end{theorem}

(Proof in Appendix~\ref{sec:appendix-carrier}.)

\begin{corollary}[Exact off-resonant branch shifts]
\label{cor:exact-off-resonant-hidden-channel-branch-shifts}\lean{WaveLane.HiddenChannel.offResonantShifts}
Assume the hypotheses of
Theorem~\ref{thm:hidden-channel-wave-remainder-and-branch-splitting}. Define
\[
\delta_{\chi}(k)
:=
\omega_{\mathrm{hid}}(k)^2-\omega_{\mathrm{vis}}(k)^2.
\]
If \(\delta_{\chi}(k)>0\) for a selected mode \(k\), then the two exact
branches satisfy
\[
\omega_-(k)^2
=
\omega_{\mathrm{vis}}(k)^2
-
\frac{2\gamma_{\chi}^2}{
\delta_{\chi}(k)+\sqrt{\delta_{\chi}(k)^2+4\gamma_{\chi}^2}
},
\]
\[
\omega_+(k)^2
=
\omega_{\mathrm{hid}}(k)^2
+
\frac{2\gamma_{\chi}^2}{
\delta_{\chi}(k)+\sqrt{\delta_{\chi}(k)^2+4\gamma_{\chi}^2}
}.
\]
In particular, away from hidden resonance the visible-like branch is shifted
downward and the hidden-like branch is shifted upward by the same exact
returned remainder amplitude.
\end{corollary}

(Proof in Appendix~\ref{sec:appendix-carrier}.)

\begin{corollary}[Exact avoided crossing of the coupled wave branches]
\label{cor:exact-avoided-crossing-coupled-wave-branches}\lean{WaveLane.HiddenChannel.avoidedCrossing}
Assume the hypotheses of
Theorem~\ref{thm:hidden-channel-wave-remainder-and-branch-splitting}. Suppose
there exists a selected mode \(k_\ast\) such that
\[
\omega_{\mathrm{vis}}(k_\ast)^2
=
\omega_{\mathrm{hid}}(k_\ast)^2
=:
\omega_\ast^2.
\]
Then
\[
\omega_{\pm}(k_\ast)^2=\omega_\ast^2\pm\gamma_{\chi},
\]
so the squared-frequency gap at the would-be crossing is exactly
\[
\omega_+(k_\ast)^2-\omega_-(k_\ast)^2=2\gamma_{\chi}.
\]
In particular, if \(\gamma_{\chi}>0\), the two branches do not cross.
\end{corollary}

(Proof in Appendix~\ref{sec:appendix-carrier}.)

\begin{corollary}[Low-energy renormalized mass and wave speed]
\label{cor:low-energy-renormalized-mass-and-wave-speed-hidden-channel}\lean{WaveLane.HiddenChannel.lowEnergyRenormalization}
Assume the hypotheses of
Theorem~\ref{thm:hidden-channel-wave-remainder-and-branch-splitting}. Define
\[
\Delta_0:=\widetilde m_{\chi}^2-m_{\chi}^2,
\qquad
\Delta_1:=\widetilde c_{\chi}^2-c_{\chi}^2,
\]
and assume \(\Delta_0>0\). Then as \(k\to 0\), the lower exact branch has the
expansion
\[
\omega_-(k)^2
=
m_{\chi,\mathrm{eff}}^2
+
c_{\chi,\mathrm{eff}}^2\,k^2
+
O(k^4),
\]
where
\[
m_{\chi,\mathrm{eff}}^2
=
\frac{
m_{\chi}^2+\widetilde m_{\chi}^2
-\sqrt{\Delta_0^2+4\gamma_{\chi}^2}
}{2},
\]
\[
c_{\chi,\mathrm{eff}}^2
=
\frac{
c_{\chi}^2+\widetilde c_{\chi}^2
-\frac{\Delta_0\Delta_1}{\sqrt{\Delta_0^2+4\gamma_{\chi}^2}}
}{2}.
\]
Thus hidden-channel return renormalizes both the visible mass gap and the
long-wave propagation speed on the selected visible-like branch.
\end{corollary}

(Proof in Appendix~\ref{sec:appendix-carrier}.)

These corollaries identify the observable content of the solved hidden-channel
wave model. The returned wave remainder shifts the visible branch away from
resonance, opens an avoided crossing at resonance, and renormalizes the
effective mass gap and long-wave propagation speed at low energy.
